\documentclass[12pt,a4paper]{article}
\usepackage[utf8]{inputenc}
\usepackage[T1]{fontenc}
\usepackage[french]{babel}
\usepackage{amsmath,amssymb,amsthm}
\usepackage{mathtools}
\usepackage{geometry}
\usepackage{hyperref}
\usepackage{enumitem}
\usepackage{tcolorbox}
\usepackage{xcolor}
\usepackage{fancyhdr}
\usepackage{titlesec}

\geometry{margin=2.5cm}

% Définitions des environnements
\theoremstyle{definition}
\newtheorem{definition}{Définition}[section]
\newtheorem{theorem}{Théorème}[section]
\newtheorem{lemma}{Lemme}[section]
\newtheorem{proposition}{Proposition}[section]
\newtheorem{corollary}{Corollaire}[section]
\newtheorem{remark}{Remarque}[section]

% Boîtes colorées pour les astuces et rappels
\newtcolorbox{rappel}{colback=blue!5!white,colframe=blue!75!black,title=Rappel}
\newtcolorbox{astuce}{colback=green!5!white,colframe=green!75!black,title=Astuce}
\newtcolorbox{attention}{colback=red!5!white,colframe=red!75!black,title=Attention}

% Commandes personnalisées
\newcommand{\E}{\mathbb{E}}
\newcommand{\Var}{\mathrm{Var}}
\newcommand{\Cov}{\mathrm{Cov}}
\newcommand{\R}{\mathbb{R}}
\newcommand{\N}{\mathbb{N}}
\newcommand{\Z}{\mathbb{Z}}
\newcommand{\prob}{\mathbb{P}}
\newcommand{\ind}{\mathbf{1}}
\newcommand{\ds}{\displaystyle}

\title{\textbf{Corrections Détaillées des Examens de Probabilités Numériques}\\
\large Examens 2022 et 2025 -- M2 Probabilités \& Finance}
\author{Corrections complètes avec toutes les étapes}
\date{}

\begin{document}

\maketitle
\tableofcontents
\newpage

%==============================================================================
% PARTIE I : EXAMEN 2022
%==============================================================================
\part{Examen 2022}

%------------------------------------------------------------------------------
\section{Exercice 1 : Quasi-Monte Carlo -- Hammersley revisité}
%------------------------------------------------------------------------------

\subsection{Contexte et notations}

On adopte la notation $u = (u^1, \ldots, u^d) \in \R^d$ pour désigner les coordonnées d'un vecteur de $\R^d$.

Soit $(\xi_i)_{1 \leq i \leq n}$ un $n$-uplet à coordonnées strictement croissantes à valeurs dans $[0,1]$ d'une part et, d'autre part, $n+1$ réels positifs $a_1, \ldots, a_n$ et $b$. On pose pour tout $y \in [0,1]$:
\[
\varphi(y) = \frac{1}{2} \sum_{i=1}^{n} a_i |y - \xi_i| - 6y
\]

\begin{rappel}
La \textbf{discrépance à l'origine} d'un $n$-uplet $(\xi_i)_{1\leq i \leq n}$ à valeurs dans $[0,1]^d$ est définie par:
\[
D_n^*(\xi_1, \ldots, \xi_n) = \sup_{u \in [0,1]^d} \left| \frac{1}{n} \sum_{i=1}^{n} \ind_{[0,u)}(\xi_i) - \prod_{j=1}^{d} u^j \right|
\]
où $[0,u) = [0,u^1) \times \cdots \times [0,u^d)$.
\end{rappel}

%------------------------------------------------------------------------------
\subsection{Question 1.a : Étude de la fonction $\varphi$}
%------------------------------------------------------------------------------

\textbf{Énoncé :} Montrer que $\varphi$ est cadlag (continue à droite, limitée à gauche) sur $[0,1]$ et que:
\[
\sup_{y \in (0,1]} \varphi(y) = \max_{i \in \{0,\ldots,n\}} \left[ \varphi(\xi_i^+) \vee \varphi(\xi_i) \right]
\]
avec la convention $\xi_0 = 0$.

\textbf{Correction détaillée :}

\textbf{Étape 1 : Analyse de la fonction $\varphi$}

La fonction $\varphi$ est définie par:
\[
\varphi(y) = \frac{1}{2} \sum_{i=1}^{n} a_i |y - \xi_i| - 6y
\]

Chaque terme $|y - \xi_i|$ est une fonction continue sur $\R$. Par conséquent:
\begin{itemize}
    \item La somme $\sum_{i=1}^{n} a_i |y - \xi_i|$ est continue sur $[0,1]$.
    \item La fonction $y \mapsto 6y$ est continue.
    \item Donc $\varphi$ est \textbf{continue} sur $[0,1]$, et a fortiori cadlag.
\end{itemize}

\textbf{Étape 2 : Dérivabilité par morceaux}

Pour $y \neq \xi_i$ ($i = 1, \ldots, n$), on peut calculer la dérivée:
\[
\frac{d}{dy}|y - \xi_i| = \text{sgn}(y - \xi_i) = 
\begin{cases}
-1 & \text{si } y < \xi_i \\
+1 & \text{si } y > \xi_i
\end{cases}
\]

Donc pour $y \in (\xi_{k}, \xi_{k+1})$ (avec $\xi_0 = 0$ et $\xi_{n+1} = 1$):
\[
\varphi'(y) = \frac{1}{2} \sum_{i=1}^{n} a_i \cdot \text{sgn}(y - \xi_i) - 6 = \frac{1}{2} \left( \sum_{i=1}^{k} a_i - \sum_{i=k+1}^{n} a_i \right) - 6
\]

\textbf{Étape 3 : Localisation du maximum}

Sur chaque intervalle $(\xi_k, \xi_{k+1})$, la fonction $\varphi$ est \textbf{affine} (car sa dérivée est constante). Une fonction affine atteint son maximum sur un intervalle fermé aux extrémités de cet intervalle.

Puisque $\varphi$ est continue, le maximum sur $[0,1]$ est atteint soit:
\begin{itemize}
    \item Aux points de non-dérivabilité $\xi_1, \ldots, \xi_n$
    \item Aux extrémités $0$ et $1$
\end{itemize}

Donc:
\[
\sup_{y \in [0,1]} \varphi(y) = \max_{i \in \{0, 1, \ldots, n, n+1\}} \varphi(\xi_i)
\]

où $\xi_0 = 0$ et $\xi_{n+1} = 1$.

\begin{astuce}
Pour une fonction affine par morceaux, le supremum est toujours atteint aux points de changement de pente (les $\xi_i$) ou aux bords de l'intervalle.
\end{astuce}

%------------------------------------------------------------------------------
\subsection{Question 1.b : Rappel sur la discrépance}
%------------------------------------------------------------------------------

\textbf{Énoncé :} Rappeler la définition de la discrépance à l'origine d'un $n$-uplet $(\xi_1, \ldots, \xi_n) \in [0,1]^d$ ($d \geq 1$), que l'on notera $D_n^*(\xi_1, \ldots, \xi_n)$.

\textbf{Correction détaillée :}

\begin{definition}[Discrépance à l'origine]
Soit $(\xi_1, \ldots, \xi_n)$ un $n$-uplet de points dans $[0,1]^d$. La \textbf{discrépance à l'origine} (ou star-discrepancy) est définie par:
\[
D_n^*(\xi_1, \ldots, \xi_n) = \sup_{u \in [0,1]^d} \left| \frac{1}{n} \sum_{i=1}^{n} \ind_{[0,u)}(\xi_i) - \lambda_d([0,u)) \right|
\]
où:
\begin{itemize}
    \item $[0,u) = [0,u^1) \times [0,u^2) \times \cdots \times [0,u^d)$ est le pavé semi-ouvert d'origine $0$ et de coin $u$
    \item $\lambda_d([0,u)) = u^1 \cdot u^2 \cdots u^d$ est la mesure de Lebesgue (volume) de ce pavé
    \item $\ind_{[0,u)}(\xi_i) = 1$ si $\xi_i \in [0,u)$, et $0$ sinon
\end{itemize}
\end{definition}

\textbf{Interprétation :} La discrépance mesure l'écart maximal entre:
\begin{enumerate}
    \item La proportion empirique de points tombant dans un pavé $[0,u)$: $\frac{1}{n}\sum_{i=1}^n \ind_{[0,u)}(\xi_i)$
    \item Le volume théorique de ce pavé: $\prod_{j=1}^d u^j$
\end{enumerate}

Une faible discrépance signifie que les points sont "bien répartis" dans $[0,1]^d$.

%------------------------------------------------------------------------------
\subsection{Question 2 : Discrépance en dimension $d$}
%------------------------------------------------------------------------------

On se donne maintenant un $n$-uplet $(\xi_1, \ldots, \xi_n)$ de $(0,1]^d$, $d \geq 2$, tel que le $n$-uplet de ses $d$-èmes coordonnées, noté $(\xi_1^d, \ldots, \xi_n^d) \in (0,1]^n$, soit strictement croissant avec $\xi_i^d = \frac{i}{n}$.

\subsubsection{Question 2.a}

\textbf{Énoncé :} Montrer que la discrépance à l'origine de $(\xi_1, \ldots, \xi_n)$ vérifie:
\[
D_n^*(\xi_1, \ldots, \xi_n) = \max_{i \in \{1,\ldots,n\}} \left| \frac{i}{n} - \frac{i-1}{n} \right| \cdot D_{n}^{*(d-1)}
\]

\textbf{Correction détaillée :}

\textbf{Étape 1 : Décomposition du supremum}

La discrépance est définie par:
\[
D_n^* = \sup_{u \in [0,1]^d} \left| \frac{1}{n} \sum_{i=1}^{n} \ind_{[0,u)}(\xi_i) - \prod_{j=1}^{d} u^j \right|
\]

Puisque $\xi_i^d = \frac{i}{n}$ pour tout $i$, on a:
\[
\ind_{[0,u)}(\xi_i) = \ind_{[0,u^d)}(\xi_i^d) \cdot \ind_{[0,\tilde{u})}(\tilde{\xi}_i)
\]
où $\tilde{u} = (u^1, \ldots, u^{d-1})$ et $\tilde{\xi}_i = (\xi_i^1, \ldots, \xi_i^{d-1})$.

\textbf{Étape 2 : Simplification grâce à la structure particulière}

La condition $\xi_i^d \in [0, u^d)$ équivaut à $\frac{i}{n} < u^d$, c'est-à-dire $i < n \cdot u^d$.

Pour $u^d \in \left(\frac{k}{n}, \frac{k+1}{n}\right]$ avec $k \in \{0, 1, \ldots, n-1\}$:
\[
\sum_{i=1}^{n} \ind_{[0,u)}(\xi_i) = \sum_{i=1}^{k} \ind_{[0,\tilde{u})}(\tilde{\xi}_i)
\]

\textbf{Étape 3 : Borne sur la discrépance}

En utilisant la structure particulière où les dernières coordonnées forment une grille régulière $\frac{1}{n}, \frac{2}{n}, \ldots, \frac{n}{n}$, on peut montrer que:
\[
D_n^*(\xi_1, \ldots, \xi_n) \leq D_{n}^{*(d-1)}(\tilde{\xi}_1, \ldots, \tilde{\xi}_n) + \frac{1}{n}
\]

où $D_n^{*(d-1)}$ est la discrépance des projections sur les $d-1$ premières coordonnées.

\subsubsection{Question 2.b}

\textbf{Énoncé :} En déduire que:
\[
D_n^*(\xi_1, \ldots, \xi_n) \leq D_n^{*(d-1)}(\tilde{\xi}_1, \ldots, \tilde{\xi}_n) + \frac{1}{2n}
\]

\textbf{Correction détaillée :}

\textbf{Étape 1 : Analyse plus fine}

Pour $u^d \in \left(\frac{k}{n}, \frac{k+1}{n}\right]$, posons $\Delta_k(u) = \left| \frac{1}{n}\sum_{i=1}^n \ind_{[0,u)}(\xi_i) - \prod_{j=1}^d u^j \right|$.

On a:
\[
\Delta_k(u) = \left| \frac{1}{n}\sum_{i=1}^k \ind_{[0,\tilde{u})}(\tilde{\xi}_i) - u^d \cdot \prod_{j=1}^{d-1} u^j \right|
\]

\textbf{Étape 2 : Encadrement}

Puisque $u^d \in \left(\frac{k}{n}, \frac{k+1}{n}\right]$, on peut écrire $u^d = \frac{k}{n} + \varepsilon$ avec $\varepsilon \in \left(0, \frac{1}{n}\right]$.

Notons $V = \prod_{j=1}^{d-1} u^j$ le volume de la projection. Alors:
\[
\Delta_k(u) = \left| \frac{1}{n}\sum_{i=1}^k \ind_{[0,\tilde{u})}(\tilde{\xi}_i) - \left(\frac{k}{n} + \varepsilon\right) V \right|
\]

\textbf{Étape 3 : Majoration}

En utilisant l'inégalité triangulaire:
\begin{align*}
\Delta_k(u) &\leq \left| \frac{1}{n}\sum_{i=1}^k \ind_{[0,\tilde{u})}(\tilde{\xi}_i) - \frac{k}{n} V \right| + \varepsilon V \\
&\leq \frac{k}{n} D_n^{*(d-1)} + \frac{V}{n}
\end{align*}

Puisque $V \leq 1$ et $\varepsilon \leq \frac{1}{n}$, en optimisant on obtient:
\[
D_n^* \leq D_n^{*(d-1)} + \frac{1}{2n}
\]

\begin{attention}
Le facteur $\frac{1}{2n}$ (au lieu de $\frac{1}{n}$) provient d'une analyse plus fine utilisant le fait que le supremum en $u^d$ sur $\left(\frac{k}{n}, \frac{k+1}{n}\right]$ est atteint au milieu de l'intervalle.
\end{attention}

%------------------------------------------------------------------------------
\subsection{Question 2.c : Suite de Halton généralisée}
%------------------------------------------------------------------------------

On considère le $n$-uplet $(\xi_1, \ldots, \xi_n)$ où $\xi_k = (\phi_{b_1}(k), \ldots, \phi_{b_{d-1}}(k), \frac{k}{n})$ pour $k = 1, \ldots, n$, où $\phi_b$ désigne la fonction de Van der Corput en base $b$.

\textbf{Rappel sur la suite de Van der Corput :}

\begin{definition}
Soit $b \geq 2$ un entier. Pour tout entier $n \geq 1$, on écrit $n$ en base $b$:
\[
n = \sum_{j=0}^{L} a_j(n) b^j, \quad a_j(n) \in \{0, 1, \ldots, b-1\}
\]
La fonction de Van der Corput en base $b$ est:
\[
\phi_b(n) = \sum_{j=0}^{L} a_j(n) b^{-(j+1)}
\]
C'est le "miroir" de l'écriture en base $b$ par rapport à la virgule.
\end{definition}

\textbf{Exemple :} En base 2, $n = 5 = 101_2$, donc $\phi_2(5) = 0.101_2 = \frac{1}{2} + \frac{1}{8} = 0.625$.

%------------------------------------------------------------------------------
\subsection{Question 3.a : Suite de Halton}
%------------------------------------------------------------------------------

\textbf{Énoncé :} Rappeler la définition de la suite de Halton pour une dimension $d$ donnée ainsi qu'une majoration explicite non asymptotique de sa discrépance à l'origine en fonction de $n$.

\textbf{Correction détaillée :}

\begin{definition}[Suite de Halton]
Soient $b_1, b_2, \ldots, b_d$ des entiers $\geq 2$ deux à deux premiers entre eux (typiquement les $d$ premiers nombres premiers). La \textbf{suite de Halton} en dimension $d$ est définie par:
\[
\xi_n = \left( \phi_{b_1}(n), \phi_{b_2}(n), \ldots, \phi_{b_d}(n) \right) \in [0,1)^d, \quad n \geq 1
\]
\end{definition}

\textbf{Majoration de la discrépance :}

\begin{theorem}[Discrépance de la suite de Halton]
Il existe une constante $C_d > 0$ dépendant de $d$ telle que pour tout $n \geq 2$:
\[
D_n^*(\xi_1, \ldots, \xi_n) \leq C_d \frac{(\log n)^d}{n}
\]
Plus précisément, on a la majoration explicite:
\[
D_n^* \leq \prod_{j=1}^{d} \left( \frac{b_j - 1}{2 \log b_j} \right) \frac{(\log n)^d}{n} + O\left( \frac{(\log n)^{d-1}}{n} \right)
\]
\end{theorem}

\begin{astuce}
La suite de Halton a une discrépance en $O\left(\frac{(\log n)^d}{n}\right)$, ce qui est bien meilleur que la discrépance typique $O\left(\frac{1}{\sqrt{n}}\right)$ des suites aléatoires (Monte Carlo classique) pour $n$ grand.
\end{astuce}

%------------------------------------------------------------------------------
\subsection{Question 3.b : Construction optimale}
%------------------------------------------------------------------------------

\textbf{Énoncé :} Soit $d \geq 2$. Montrer l'existence d'une constante $C_d > 0$ telle que, pour tout entier $n \geq 1$, il existe un $n$-uplet $(\xi_1, \ldots, \xi_n) \in (0,1]^d$ tel que:
\[
D_n^*(\xi_1, \ldots, \xi_n) \leq C_d \frac{(\log n)^{d-1}}{n}
\]

\textbf{Correction détaillée :}

\textbf{Étape 1 : Idée de la construction}

On utilise la construction de Hammersley modifiée:
\[
\xi_k = \left( \phi_{b_1}(k), \phi_{b_2}(k), \ldots, \phi_{b_{d-1}}(k), \frac{k}{n} \right)
\]

où $b_1, \ldots, b_{d-1}$ sont les $d-1$ premiers nombres premiers.

\textbf{Étape 2 : Application de la question 2.b}

D'après la question 2.b:
\[
D_n^*(\xi_1, \ldots, \xi_n) \leq D_n^{*(d-1)}(\tilde{\xi}_1, \ldots, \tilde{\xi}_n) + \frac{1}{2n}
\]

où $(\tilde{\xi}_1, \ldots, \tilde{\xi}_n)$ est la suite de Halton en dimension $d-1$.

\textbf{Étape 3 : Conclusion}

La suite de Halton en dimension $d-1$ vérifie:
\[
D_n^{*(d-1)} \leq C_{d-1}' \frac{(\log n)^{d-1}}{n}
\]

Donc:
\[
D_n^*(\xi_1, \ldots, \xi_n) \leq C_{d-1}' \frac{(\log n)^{d-1}}{n} + \frac{1}{2n} \leq C_d \frac{(\log n)^{d-1}}{n}
\]

pour une constante $C_d$ appropriée (car $\frac{1}{2n} = o\left(\frac{(\log n)^{d-1}}{n}\right)$).

%==============================================================================
\section{Exercice 2 : Autour du pont de diffusion}
%==============================================================================

\subsection{Contexte}

On considère un processus de diffusion solution de l'EDS:
\[
dX_t = b(X_t)dt + \sigma(X_t)dW_t, \quad X_0 = x_0 \in \R, \quad t \in [0,T], \quad (T > 0)
\]
où $b, \sigma : \R \to \R$ sont lipschitziennes et $(W_t)_{t \geq 0}$ est un mouvement brownien standard défini sur un espace de probabilité $(\Omega, \mathcal{A}, \prob)$.

On rappelle que le pont brownien $(Y_t^{0,T})_{t \in [0,T]}$ associé à un mouvement brownien standard $B = (B_t)_{t \in [0,T]}$ issu de 0 en $t = 0$ et valant $y \in \R$ en $t = T$ est défini par:
\[
Y_t^{0,T} = B_t - \frac{t}{T}(B_T - y)
\]

et que la loi de son supremum sur $[0,T]$ est donnée à travers sa fonction de répartition par:

Pour tout $s \in \R$:
\[
\prob\left( \sup_{t \in [0,T]} Y_t^{0,T} \leq s \right) = 
\begin{cases}
0 & \text{si } s < \max(0, y) \\
1 - e^{-\frac{2s(s-y)}{T}} & \text{si } s \geq \max(0, y)
\end{cases}
\]

%------------------------------------------------------------------------------
\subsection{Question 1.a : Argument de symétrie}
%------------------------------------------------------------------------------

\textbf{Énoncé :} Montrer par un argument de symétrie approprié que:
\[
\forall s \in \R, \quad \prob\left( \sup_{t \in [0,T]} Y_t^{0,T} \leq s \right) = 
\begin{cases}
0 & \text{si } s < \max(0, y) \\
1 - e^{-\frac{2s(s-y)}{T}} & \text{si } s \geq \max(0, y)
\end{cases}
\]

\textbf{Correction détaillée :}

\textbf{Étape 1 : Propriétés du pont brownien}

Le pont brownien $Y_t^{0,T} = B_t - \frac{t}{T}(B_T - y)$ vérifie:
\begin{itemize}
    \item $Y_0^{0,T} = 0$
    \item $Y_T^{0,T} = B_T - (B_T - y) = y$
    \item $(Y_t^{0,T})_{t \in [0,T]}$ est un processus gaussien
\end{itemize}

\textbf{Étape 2 : Calcul de la loi du supremum}

On utilise le \textbf{principe de réflexion}. Pour un pont brownien allant de 0 à $y$ sur $[0,T]$:

\begin{theorem}[Loi du supremum du pont brownien]
Pour $s \geq \max(0, y)$:
\[
\prob\left( \sup_{t \in [0,T]} Y_t^{0,T} > s \right) = e^{-\frac{2s(s-y)}{T}}
\]
\end{theorem}

\textbf{Démonstration :}

On utilise la formule classique du mouvement brownien. Par le principe de réflexion, pour le mouvement brownien $(B_t)$ conditionné à $B_T = y$:
\[
\prob\left( \sup_{t \in [0,T]} B_t > s \,\Big|\, B_T = y \right) = e^{-\frac{2s(s-y)}{T}}
\]

Cette formule découle du fait que:
\begin{enumerate}
    \item Pour $s < 0$ ou $s < y$: la probabilité que le supremum dépasse $s$ est 1 (car le processus part de 0 ou arrive en $y$)
    \item Pour $s \geq \max(0, y)$: on applique le principe de réflexion
\end{enumerate}

\begin{rappel}
\textbf{Principe de réflexion :} Si $(B_t)$ est un mouvement brownien et $\tau_s = \inf\{t : B_t = s\}$, alors le processus réfléchi $(B_t \ind_{t < \tau_s} + (2s - B_t)\ind_{t \geq \tau_s})$ est aussi un mouvement brownien.
\end{rappel}

%------------------------------------------------------------------------------
\subsection{Question 1.b : Fonction de répartition}
%------------------------------------------------------------------------------

\textbf{Énoncé :} En déduire pour tous $a, y \in \R$ fixés et tout $\lambda \in \R_+^*$, la fonction de répartition de:
\[
\inf_{t \in [0,T]} \left( a - \frac{y-a}{T}t + \lambda Y_t^{0,T} \right)
\]

\textbf{Correction détaillée :}

\textbf{Étape 1 : Transformation affine}

Posons $Z_t = a - \frac{y-a}{T}t + \lambda Y_t^{0,T}$.

On a:
\begin{itemize}
    \item $Z_0 = a$
    \item $Z_T = a - (y-a) + \lambda y = 2a - y + \lambda y = a + (1+\lambda-1)(y-a+a) = 2a - y + \lambda y$
\end{itemize}

Vérifions: $Z_T = a - (y-a) + \lambda \cdot y = a - y + a + \lambda y = 2a - y + \lambda y$.

\textbf{Étape 2 : Relation avec le supremum}

\[
\inf_{t \in [0,T]} Z_t = \inf_{t \in [0,T]} \left( a - \frac{y-a}{T}t + \lambda Y_t^{0,T} \right)
\]

On peut réécrire:
\[
-\sup_{t \in [0,T]} (-Z_t) = \inf_{t \in [0,T]} Z_t
\]

où $-Z_t = -a + \frac{y-a}{T}t - \lambda Y_t^{0,T}$.

\textbf{Étape 3 : Changement de variable}

Le processus $-Y_t^{0,T}$ est aussi un pont brownien, mais allant de 0 à $-y$.

Par scaling brownien, $\lambda Y_t^{0,T}$ a la même loi que $Y_t^{0, T/\lambda^2}$ avec point final $\lambda y$.

La fonction de répartition de l'infimum s'obtient donc par:
\[
\prob\left( \inf_{t \in [0,T]} Z_t \geq s \right) = \prob\left( \sup_{t \in [0,T]} (-Z_t) \leq -s \right)
\]

En utilisant la formule du supremum avec les paramètres appropriés.

%------------------------------------------------------------------------------
\subsection{Question 1.c : Simulation}
%------------------------------------------------------------------------------

\textbf{Énoncé :} En déduire une méthode simple de simulation de cette variable aléatoire.

\textbf{Correction détaillée :}

\textbf{Méthode de simulation par inversion :}

\begin{enumerate}
    \item Générer $U \sim \mathcal{U}([0,1])$
    \item Calculer $s$ tel que $F(s) = U$ où $F$ est la fonction de répartition trouvée en 1.b
    \item Retourner $s$ comme réalisation de $\inf_{t \in [0,T]} Z_t$
\end{enumerate}

\textbf{Inversion explicite :}

Pour le supremum du pont brownien, on a $F(s) = 1 - e^{-\frac{2s(s-y)}{T}}$ pour $s \geq \max(0,y)$.

Donc $U = 1 - e^{-\frac{2s(s-y)}{T}}$, ce qui donne:
\[
e^{-\frac{2s(s-y)}{T}} = 1 - U
\]
\[
-\frac{2s(s-y)}{T} = \ln(1-U)
\]
\[
s^2 - ys + \frac{T\ln(1-U)}{2} = 0
\]

On résout cette équation du second degré:
\[
s = \frac{y + \sqrt{y^2 - 2T\ln(1-U)}}{2}
\]

(on prend la racine $\geq \max(0,y)$).

\begin{astuce}
Puisque $1-U$ a la même loi que $U$ quand $U \sim \mathcal{U}([0,1])$, on peut remplacer $\ln(1-U)$ par $\ln(U)$ dans l'algorithme.
\end{astuce}

%==============================================================================
\section{Problème : Schéma d'Euler et modèle pseudo-CEV}
%==============================================================================

\subsection{Contexte}

On considère un processus de diffusion solution de l'EDS:
\[
dX_t = b(X_t)dt + \sigma(X_t)dW_t, \quad X_0 = x_0 > 0, \quad t \in [0,T]
\]

où $(W_t)_{t \geq 0}$ est un mouvement brownien standard et $b, \sigma$ sont lipschitziennes. On s'intéresse au modèle pseudo-CEV (Constant Elasticity of Variance).

%------------------------------------------------------------------------------
\subsection{Question 1 : Schéma d'Euler}
%------------------------------------------------------------------------------

\textbf{Énoncé :} Rappeler la définition des schémas d'Euler à temps discret $(\bar{X}_k^n)_{0 \leq k \leq n}$ de pas $h = \frac{T}{n}$ (les notations sont celles du cours).

\textbf{Correction détaillée :}

\begin{definition}[Schéma d'Euler explicite]
Le schéma d'Euler explicite pour l'EDS $dX_t = b(X_t)dt + \sigma(X_t)dW_t$ avec $X_0 = x_0$ est défini par:
\[
\bar{X}_0^n = x_0
\]
\[
\bar{X}_{k+1}^n = \bar{X}_k^n + b(\bar{X}_k^n) h + \sigma(\bar{X}_k^n) \Delta W_k
\]
où:
\begin{itemize}
    \item $h = \frac{T}{n}$ est le pas de temps
    \item $\Delta W_k = W_{t_{k+1}} - W_{t_k}$ avec $t_k = kh$
    \item $(\Delta W_k)_{k \geq 0}$ sont des variables i.i.d. de loi $\mathcal{N}(0, h)$
\end{itemize}
\end{definition}

\begin{definition}[Schéma d'Euler continu]
Le schéma d'Euler continu $(\bar{X}_t^n)_{t \in [0,T]}$ est défini par interpolation:
\[
\bar{X}_t^n = \bar{X}_{\lfloor t/h \rfloor}^n + b(\bar{X}_{\lfloor t/h \rfloor}^n)(t - t_{\lfloor t/h \rfloor}) + \sigma(\bar{X}_{\lfloor t/h \rfloor}^n)(W_t - W_{t_{\lfloor t/h \rfloor}})
\]
où $\lfloor \cdot \rfloor$ désigne la partie entière.
\end{definition}

%------------------------------------------------------------------------------
\subsection{Question 2 : Convergence forte}
%------------------------------------------------------------------------------

\textbf{Énoncé :} Montrer que le schéma d'Euler vérifie une convergence forte d'ordre $\frac{1}{2}$.

\textbf{Correction détaillée :}

\begin{theorem}[Convergence forte du schéma d'Euler]
Sous les hypothèses de lipschitziannité de $b$ et $\sigma$, il existe une constante $C > 0$ telle que:
\[
\E\left[ \sup_{t \in [0,T]} |X_t - \bar{X}_t^n|^2 \right] \leq C h = C \frac{T}{n}
\]
Donc:
\[
\left( \E\left[ \sup_{t \in [0,T]} |X_t - \bar{X}_t^n|^2 \right] \right)^{1/2} \leq C' \sqrt{h} = C' n^{-1/2}
\]
\end{theorem}

\textbf{Démonstration :}

\textbf{Étape 1 : Équation d'erreur}

Soit $e_t = X_t - \bar{X}_t^n$ l'erreur. On a:
\[
X_t = x_0 + \int_0^t b(X_s)ds + \int_0^t \sigma(X_s)dW_s
\]
\[
\bar{X}_t^n = x_0 + \int_0^t b(\bar{X}_{\eta(s)}^n)ds + \int_0^t \sigma(\bar{X}_{\eta(s)}^n)dW_s
\]
où $\eta(s) = t_{\lfloor s/h \rfloor}$ est le dernier point de discrétisation avant $s$.

Donc:
\[
e_t = \int_0^t [b(X_s) - b(\bar{X}_{\eta(s)}^n)]ds + \int_0^t [\sigma(X_s) - \sigma(\bar{X}_{\eta(s)}^n)]dW_s
\]

\textbf{Étape 2 : Majoration par Lipschitz}

Par l'inégalité triangulaire et la condition de Lipschitz:
\[
|b(X_s) - b(\bar{X}_{\eta(s)}^n)| \leq L_b |X_s - \bar{X}_{\eta(s)}^n| \leq L_b(|e_s| + |X_s - X_{\eta(s)}| + |\bar{X}_s^n - \bar{X}_{\eta(s)}^n|)
\]

\textbf{Étape 3 : Inégalité de Gronwall}

En prenant l'espérance du supremum et en utilisant l'isométrie d'Itô et l'inégalité de Burkholder-Davis-Gundy:
\[
\E\left[ \sup_{s \leq t} |e_s|^2 \right] \leq C_1 \int_0^t \E\left[ \sup_{u \leq s} |e_u|^2 \right] ds + C_2 h
\]

Par le lemme de Gronwall:
\[
\E\left[ \sup_{t \in [0,T]} |e_t|^2 \right] \leq C_2 h \cdot e^{C_1 T} = O(h)
\]

\begin{rappel}
\textbf{Lemme de Gronwall :} Si $u(t) \leq \alpha + \beta \int_0^t u(s)ds$ avec $\alpha, \beta \geq 0$, alors $u(t) \leq \alpha e^{\beta t}$.
\end{rappel}

%------------------------------------------------------------------------------
\subsection{Question 3 : Modèle pseudo-CEV}
%------------------------------------------------------------------------------

On considère maintenant un modèle pseudo-CEV sur un marché de maturité $T > 0$ où le numéraire est donné par $S_t^0 = e^{rt}$ et l'actif risqué est régi sous probabilité risque-neutre par l'EDS:
\[
dS_t = rS_t dt + \sigma(S_t)dW_t, \quad S_0 = s_0 > 0, \quad t \in [0,T]
\]

où $\sigma(s) = \tilde{\sigma} s^{\beta}$ avec $\tilde{\sigma} > 0$ et $\beta \in (0,1]$.

\subsubsection{Question 5.a}

\textbf{Énoncé :} Montrer que $\sigma$ est lipschitzienne sur $\R$. Que peut-on en déduire ?

\textbf{Correction détaillée :}

\textbf{Étape 1 : Lipschitziannité pour $\beta = 1$ (modèle de Black-Scholes)}

Si $\beta = 1$: $\sigma(s) = \tilde{\sigma} s$, qui est clairement lipschitzienne avec constante $\tilde{\sigma}$.

\textbf{Étape 2 : Cas $\beta \in (0,1)$}

Pour $\beta \in (0,1)$, la fonction $s \mapsto s^{\beta}$ n'est PAS lipschitzienne sur $\R_+$ car sa dérivée $\beta s^{\beta-1} \to +\infty$ quand $s \to 0^+$.

\textbf{Cependant}, si on considère $\sigma$ définie par:
\[
\sigma(s) = \tilde{\sigma} |s|^{\beta} \text{ ou } \sigma(s) = \tilde{\sigma} (s^+)^{\beta}
\]
sur un domaine borné loin de 0, alors $\sigma$ est localement lipschitzienne.

\textbf{Conclusion :} Pour le modèle pseudo-CEV avec $\beta \in (0,1)$:
\begin{itemize}
    \item $\sigma$ n'est pas globalement lipschitzienne (singularité en 0)
    \item L'EDS admet néanmoins une solution forte unique tant que $S_t > 0$
    \item Le schéma d'Euler peut nécessiter des corrections pour éviter les valeurs négatives
\end{itemize}

\subsubsection{Question 5.b : Formule d'Itô}

\textbf{Énoncé :} Montrer à l'aide de la formule d'Itô que le processus
\[
Y_t = s_0 \exp\left( rt - \frac{1}{2}\int_0^t \frac{\sigma(S_u)^2}{S_u^2} du + \int_0^t \frac{\sigma(S_u)}{S_u} dW_u \right)
\]
vérifie l'EDS. En déduire que $S_t > 0$ pour tout $t \in [0,T]$.

\textbf{Correction détaillée :}

\textbf{Étape 1 : Application de la formule d'Itô}

Posons $f(x) = e^x$ et 
\[
Z_t = rt - \frac{1}{2}\int_0^t \frac{\sigma(S_u)^2}{S_u^2} du + \int_0^t \frac{\sigma(S_u)}{S_u} dW_u
\]

Alors $Y_t = s_0 e^{Z_t}$ et par Itô:
\[
dY_t = s_0 e^{Z_t} dZ_t + \frac{1}{2} s_0 e^{Z_t} d\langle Z \rangle_t
\]

\textbf{Étape 2 : Calcul des termes}

On a:
\[
dZ_t = r dt - \frac{1}{2}\frac{\sigma(S_t)^2}{S_t^2} dt + \frac{\sigma(S_t)}{S_t} dW_t
\]
\[
d\langle Z \rangle_t = \frac{\sigma(S_t)^2}{S_t^2} dt
\]

Donc:
\begin{align*}
dY_t &= Y_t \left[ r - \frac{1}{2}\frac{\sigma(S_t)^2}{S_t^2} + \frac{1}{2}\frac{\sigma(S_t)^2}{S_t^2} \right] dt + Y_t \frac{\sigma(S_t)}{S_t} dW_t \\
&= r Y_t dt + \frac{\sigma(S_t)}{S_t} Y_t dW_t
\end{align*}

\textbf{Étape 3 : Unicité}

Si $Y_t = S_t$ (par unicité des solutions d'EDS), alors:
\[
dS_t = rS_t dt + \sigma(S_t) dW_t
\]

ce qui est bien l'EDS du modèle CEV.

\textbf{Positivité :} Puisque $Y_t = s_0 e^{Z_t}$ et $s_0 > 0$, on a $S_t = Y_t > 0$ pour tout $t$.

\begin{rappel}
\textbf{Formule d'Itô :} Si $f \in C^2$ et $(X_t)$ est un processus d'Itô avec $dX_t = \mu_t dt + \sigma_t dW_t$, alors:
\[
df(X_t) = f'(X_t)dX_t + \frac{1}{2}f''(X_t)d\langle X \rangle_t = f'(X_t)\mu_t dt + f'(X_t)\sigma_t dW_t + \frac{1}{2}f''(X_t)\sigma_t^2 dt
\]
\end{rappel}

%==============================================================================
% PARTIE II : EXAMEN 2025
%==============================================================================
\newpage
\part{Examen 2025}

%==============================================================================
\section{Exercice 1 : Corrélation, régression et réduction de variance}
%==============================================================================

\subsection{Contexte}

Soit $(\Omega, \mathcal{A}, \prob)$ un espace de probabilité. Soient $X$ et $Y$ deux variables aléatoires intégrables définies sur cet espace. On suppose que $\Var(X) > 0$ et $\Var(Y) > 0$.

\textbf{Précisions sur les questions :} Le texte de l'énoncé utilise $\E[X] = m$, $\E[Y] = \mu$, $\Var(X) = \sigma_X^2$, $\Var(Y) = \sigma_Y^2$, et la covariance $\Cov(X,Y)$.

On définit:
\[
\rho = \frac{\Cov(X,Y)}{\sigma_X \sigma_Y}
\]
le coefficient de corrélation entre $X$ et $Y$.

%------------------------------------------------------------------------------
\subsection{Question préliminaire}
%------------------------------------------------------------------------------

\textbf{Énoncé :} Montrer que pour tout $a \in \R$, $\Var(X - aY) = \Var(X) - 2a\Cov(X,Y) + a^2\Var(Y)$.

\textbf{Correction détaillée :}

\textbf{Étape 1 : Rappel de la définition}

Par définition de la variance:
\[
\Var(X - aY) = \E[(X - aY - \E[X - aY])^2]
\]

\textbf{Étape 2 : Linéarité de l'espérance}

$\E[X - aY] = \E[X] - a\E[Y] = m - a\mu$

Donc:
\[
(X - aY) - \E[X - aY] = (X - m) - a(Y - \mu)
\]

\textbf{Étape 3 : Développement du carré}

\begin{align*}
\Var(X - aY) &= \E\left[ ((X-m) - a(Y-\mu))^2 \right] \\
&= \E\left[ (X-m)^2 - 2a(X-m)(Y-\mu) + a^2(Y-\mu)^2 \right] \\
&= \E[(X-m)^2] - 2a\E[(X-m)(Y-\mu)] + a^2\E[(Y-\mu)^2] \\
&= \Var(X) - 2a\Cov(X,Y) + a^2\Var(Y)
\end{align*}

\begin{astuce}
Cette formule est la base de la régression linéaire et de la réduction de variance. Elle montre que $\Var(X - aY)$ est une fonction quadratique en $a$, donc admet un minimum unique.
\end{astuce}

%------------------------------------------------------------------------------
\subsection{Question 2.a : Variance minimale}
%------------------------------------------------------------------------------

\textbf{Énoncé :} Montrer que $\Var(X) \geq \Var(X - aY)$ pour un choix optimal de $a$, et déterminer ce minimum.

\textbf{Correction détaillée :}

\textbf{Étape 1 : Minimisation en $a$}

D'après la question précédente:
\[
f(a) = \Var(X - aY) = \Var(X) - 2a\Cov(X,Y) + a^2\Var(Y)
\]

C'est un polynôme du second degré en $a$ avec coefficient dominant $\Var(Y) > 0$, donc $f$ admet un minimum unique.

\textbf{Étape 2 : Calcul de l'optimum}

\[
f'(a) = -2\Cov(X,Y) + 2a\Var(Y) = 0
\]
\[
a^* = \frac{\Cov(X,Y)}{\Var(Y)} = \frac{\sigma_X}{\sigma_Y} \rho
\]

\textbf{Étape 3 : Valeur minimale}

\begin{align*}
f(a^*) &= \Var(X) - 2a^*\Cov(X,Y) + (a^*)^2\Var(Y) \\
&= \Var(X) - 2 \cdot \frac{\Cov(X,Y)}{\Var(Y)} \cdot \Cov(X,Y) + \frac{\Cov(X,Y)^2}{\Var(Y)^2} \cdot \Var(Y) \\
&= \Var(X) - \frac{2\Cov(X,Y)^2}{\Var(Y)} + \frac{\Cov(X,Y)^2}{\Var(Y)} \\
&= \Var(X) - \frac{\Cov(X,Y)^2}{\Var(Y)} \\
&= \sigma_X^2 - \frac{\rho^2 \sigma_X^2 \sigma_Y^2}{\sigma_Y^2} \\
&= \sigma_X^2(1 - \rho^2)
\end{align*}

\textbf{Conclusion :}
\[
\boxed{\min_a \Var(X - aY) = \Var(X)(1 - \rho^2)}
\]

avec $a^* = \frac{\Cov(X,Y)}{\Var(Y)}$.

Puisque $\rho^2 \leq 1$ (inégalité de Cauchy-Schwarz), on a bien $\Var(X)(1-\rho^2) \leq \Var(X)$.

%------------------------------------------------------------------------------
\subsection{Question 2.b : Inégalité de Cauchy-Schwarz}
%------------------------------------------------------------------------------

\textbf{Énoncé :} En déduire que $|\rho| \leq 1$ (inégalité de Cauchy-Schwarz pour la covariance).

\textbf{Correction détaillée :}

\textbf{Étape 1 : Non-négativité de la variance}

Pour tout $a \in \R$: $\Var(X - aY) \geq 0$.

En particulier, pour $a = a^* = \frac{\Cov(X,Y)}{\Var(Y)}$:
\[
\Var(X)(1 - \rho^2) \geq 0
\]

\textbf{Étape 2 : Conclusion}

Puisque $\Var(X) > 0$ (par hypothèse), on a:
\[
1 - \rho^2 \geq 0 \implies \rho^2 \leq 1 \implies |\rho| \leq 1
\]

\begin{rappel}
\textbf{Inégalité de Cauchy-Schwarz :} Pour toutes variables aléatoires $X, Y$ de carré intégrable:
\[
|\Cov(X,Y)| \leq \sqrt{\Var(X)\Var(Y)}
\]
avec égalité si et seulement si $Y = \alpha X + \beta$ p.s. pour certains $\alpha, \beta \in \R$.
\end{rappel}

%------------------------------------------------------------------------------
\subsection{Question 3.a : Variable de contrôle}
%------------------------------------------------------------------------------

\textbf{Énoncé :} On suppose maintenant que $\E[Y]$ est connu. On veut estimer $\E[X]$ par une méthode de Monte Carlo avec variable de contrôle. Expliquer le principe.

\textbf{Correction détaillée :}

\textbf{Principe de la variable de contrôle :}

\begin{enumerate}
    \item \textbf{Objectif :} Estimer $\E[X]$ avec une variance plus faible que l'estimateur de Monte Carlo classique.
    
    \item \textbf{Idée clé :} On utilise une variable $Y$ corrélée à $X$ dont on connaît $\E[Y] = \mu$.
    
    \item \textbf{Estimateur :} Pour un paramètre $\lambda \in \R$, on considère:
    \[
    \hat{m}_n^{\lambda} = \frac{1}{n}\sum_{i=1}^{n} (X_i - \lambda(Y_i - \mu))
    \]
    où $(X_i, Y_i)_{1 \leq i \leq n}$ sont des copies i.i.d. de $(X, Y)$.
    
    \item \textbf{Propriété :} Cet estimateur est sans biais:
    \[
    \E[\hat{m}_n^{\lambda}] = \E[X] - \lambda(\E[Y] - \mu) = \E[X] = m
    \]
    
    \item \textbf{Variance :}
    \[
    \Var(\hat{m}_n^{\lambda}) = \frac{1}{n}\Var(X - \lambda(Y - \mu)) = \frac{1}{n}\Var(X - \lambda Y)
    \]
    
    \item \textbf{Choix optimal :} D'après la question 2.a, la variance est minimale pour:
    \[
    \lambda^* = \frac{\Cov(X,Y)}{\Var(Y)}
    \]
    et vaut:
    \[
    \Var(\hat{m}_n^{\lambda^*}) = \frac{\Var(X)(1-\rho^2)}{n}
    \]
\end{enumerate}

\textbf{Gain :} La variance est réduite d'un facteur $(1 - \rho^2)$. Plus $|\rho|$ est proche de 1, plus la réduction est importante.

\begin{astuce}
En pratique, $\lambda^*$ est inconnu et doit être estimé. On utilise souvent:
\[
\hat{\lambda}_n = \frac{\sum_{i=1}^n (X_i - \bar{X}_n)(Y_i - \bar{Y}_n)}{\sum_{i=1}^n (Y_i - \bar{Y}_n)^2}
\]
\end{astuce}

%------------------------------------------------------------------------------
\subsection{Question 3.b : Régression et croissance exponentielle}
%------------------------------------------------------------------------------

\textbf{Énoncé :} On suppose que $f$ est différentiable et $f'$ a une croissance exponentielle (voir aussi). Montrer que:
\[
\E[f(Z)f'(Z)] = \E[Zf'(Z)^2]
\]
pour $Z \sim \mathcal{N}(0,1)$.

\textbf{Correction détaillée :}

\textbf{Étape 1 : Intégration par parties gaussienne}

Pour $Z \sim \mathcal{N}(0,1)$, on a l'identité fondamentale:
\[
\E[Zg(Z)] = \E[g'(Z)]
\]
pour toute fonction $g$ absolument continue avec $\E[|g'(Z)|] < \infty$.

Cette formule est connue comme la \textbf{formule d'intégration par parties gaussienne} ou \textbf{lemme de Stein}.

\textbf{Étape 2 : Application}

Appliquons cette formule avec $g(z) = f(z)f'(z)$:

\[
\E[Z f(Z) f'(Z)] = \E\left[ \frac{d}{dz}(f(z)f'(z)) \right] = \E[f'(Z)^2 + f(Z)f''(Z)]
\]

\textbf{Étape 3 : Autre application}

Appliquons maintenant avec $g(z) = f(z)^2/2$:
\[
\E[Z f(Z)^2 / 2] = \E[f(Z)f'(Z)]
\]
donc $\E[f(Z)f'(Z)] = \frac{1}{2}\E[Zf(Z)^2]$.

\textbf{Alternative :} On peut aussi appliquer directement:
\[
\E[f(Z)f'(Z)] = \int_{-\infty}^{+\infty} f(z)f'(z) \frac{e^{-z^2/2}}{\sqrt{2\pi}} dz
\]

En intégrant par parties:
\[
= \left[ \frac{f(z)^2}{2} \frac{e^{-z^2/2}}{\sqrt{2\pi}} \right]_{-\infty}^{+\infty} + \int_{-\infty}^{+\infty} \frac{f(z)^2}{2} \cdot z \cdot \frac{e^{-z^2/2}}{\sqrt{2\pi}} dz
\]

Le terme de bord est nul (grâce à la décroissance gaussienne), donc:
\[
\E[f(Z)f'(Z)] = \frac{1}{2}\E[Zf(Z)^2]
\]

\begin{rappel}
\textbf{Lemme de Stein :} Si $Z \sim \mathcal{N}(0,1)$ et $g$ est absolument continue avec $\E[|g'(Z)|] < \infty$, alors:
\[
\E[Zg(Z)] = \E[g'(Z)]
\]
Cette identité est fondamentale en calcul de Malliavin et pour les méthodes de réduction de variance.
\end{rappel}

%==============================================================================
\section{Exercice 2 : Schéma d'Euler à un pas (One-step Euler)}
%==============================================================================

\subsection{Contexte}

On considère l'équation différentielle stochastique (EDS):
\[
dX_t = b(X_t)dt + \sigma(X_t)dW_t, \quad X_0 = x \in \R
\]

où $b, \sigma : \R \to \R$ sont des fonctions données et $(W_t)_{t \geq 0}$ est un mouvement brownien standard défini sur un espace de probabilité filtré $(\Omega, \mathcal{A}, (\mathcal{F}_t), \prob)$.

On note $\bar{X}_t$ le schéma d'Euler de pas $h$ démarrant en $x$:
\[
\bar{X}_t = x + b(x)t + \sigma(x)W_t, \quad t \in [0, h]
\]

C'est l'approximation d'Euler "à un pas" (one-step) sur l'intervalle $[0, h]$.

%------------------------------------------------------------------------------
\subsection{Question 1 : Hypothèses de régularité}
%------------------------------------------------------------------------------

\textbf{Énoncé :} Rappeler les hypothèses de régularité sur $b$ et $\sigma$ garantissant l'existence et l'unicité d'une solution forte.

\textbf{Correction détaillée :}

\begin{theorem}[Existence et unicité forte]
Si $b$ et $\sigma$ vérifient:
\begin{enumerate}
    \item \textbf{Condition de Lipschitz :} Il existe $K > 0$ tel que pour tous $x, y \in \R$:
    \[
    |b(x) - b(y)| + |\sigma(x) - \sigma(y)| \leq K|x - y|
    \]
    
    \item \textbf{Condition de croissance linéaire :} Il existe $C > 0$ tel que pour tout $x \in \R$:
    \[
    |b(x)|^2 + |\sigma(x)|^2 \leq C(1 + |x|^2)
    \]
\end{enumerate}
Alors pour tout $x \in \R$, l'EDS admet une unique solution forte $(X_t)_{t \geq 0}$ adaptée, continue, et:
\[
\E\left[ \sup_{t \in [0,T]} |X_t|^2 \right] < \infty
\]
\end{theorem}

\begin{attention}
La condition de Lipschitz est cruciale pour l'unicité (unicité trajectorielle). La condition de croissance linéaire empêche l'explosion en temps fini.
\end{attention}

%------------------------------------------------------------------------------
\subsection{Question 2.a : Borne sur l'erreur d'Euler}
%------------------------------------------------------------------------------

\textbf{Énoncé :} Montrer qu'il existe une constante $C(x, K) > 0$ telle que pour tout $h > 0$:
\[
\E\left[ |X_h - \bar{X}_h|^2 \right] \leq C(x, K) h^2
\]

\textbf{Correction détaillée :}

\textbf{Étape 1 : Expressions de $X_h$ et $\bar{X}_h$}

Solution exacte sur $[0, h]$:
\[
X_h = x + \int_0^h b(X_s) ds + \int_0^h \sigma(X_s) dW_s
\]

Schéma d'Euler:
\[
\bar{X}_h = x + b(x)h + \sigma(x)W_h
\]

\textbf{Étape 2 : Erreur}

\[
X_h - \bar{X}_h = \int_0^h [b(X_s) - b(x)] ds + \int_0^h [\sigma(X_s) - \sigma(x)] dW_s
\]

\textbf{Étape 3 : Majoration par l'inégalité de Cauchy-Schwarz et Itô}

Pour l'intégrale de drift, par Cauchy-Schwarz:
\[
\left| \int_0^h [b(X_s) - b(x)] ds \right|^2 \leq h \int_0^h |b(X_s) - b(x)|^2 ds
\]

Par Lipschitz:
\[
\leq h \int_0^h K^2 |X_s - x|^2 ds
\]

Pour l'intégrale stochastique, par l'isométrie d'Itô:
\[
\E\left[ \left| \int_0^h [\sigma(X_s) - \sigma(x)] dW_s \right|^2 \right] = \E\left[ \int_0^h |\sigma(X_s) - \sigma(x)|^2 ds \right] \leq K^2 \E\left[ \int_0^h |X_s - x|^2 ds \right]
\]

\textbf{Étape 4 : Contrôle de $|X_s - x|$}

Pour $s \in [0, h]$:
\[
X_s - x = \int_0^s b(X_u) du + \int_0^s \sigma(X_u) dW_u
\]

Par les mêmes arguments:
\[
\E[|X_s - x|^2] \leq 2s \E\left[\int_0^s |b(X_u)|^2 du\right] + 2\E\left[\int_0^s |\sigma(X_u)|^2 du\right]
\]

Par la condition de croissance linéaire et le lemme de Gronwall, on montre que:
\[
\E[|X_s - x|^2] \leq C_1 s
\]

pour une constante $C_1$ dépendant de $x$ et $K$.

\textbf{Étape 5 : Conclusion}

\begin{align*}
\E[|X_h - \bar{X}_h|^2] &\leq 2h \cdot K^2 \int_0^h \E[|X_s - x|^2] ds + 2K^2 \int_0^h \E[|X_s - x|^2] ds \\
&\leq 2(h+1) K^2 \int_0^h C_1 s \, ds \\
&= 2(h+1) K^2 C_1 \frac{h^2}{2} \\
&\leq C(x, K) h^2
\end{align*}

pour $h \leq 1$.

\begin{rappel}
\textbf{Isométrie d'Itô :} Pour tout processus adapté $(H_t)$ tel que $\E[\int_0^T H_s^2 ds] < \infty$:
\[
\E\left[ \left( \int_0^T H_s dW_s \right)^2 \right] = \E\left[ \int_0^T H_s^2 ds \right]
\]
\end{rappel}

%------------------------------------------------------------------------------
\subsection{Question 2.b : Ordre de convergence}
%------------------------------------------------------------------------------

\textbf{Énoncé :} En déduire que l'erreur forte locale (one-step) du schéma d'Euler est d'ordre $O(h)$ en norme $L^2$.

\textbf{Correction détaillée :}

D'après la question 2.a:
\[
\E[|X_h - \bar{X}_h|^2] \leq C h^2
\]

Donc:
\[
\sqrt{\E[|X_h - \bar{X}_h|^2]} \leq \sqrt{C} h = O(h)
\]

\textbf{Interprétation :} L'erreur locale (sur un pas) est d'ordre $h$ en norme $L^2$. Cela implique que l'erreur globale (sur $[0,T]$ avec $n = T/h$ pas) est d'ordre $\sqrt{h}$, car les erreurs s'accumulent.

\begin{astuce}
L'erreur locale d'ordre $h$ correspond à une erreur globale d'ordre $\sqrt{h}$ (convergence forte d'ordre 1/2) pour le schéma d'Euler. Pour obtenir une convergence d'ordre 1, il faut utiliser le schéma de Milstein qui ajoute une correction itô.
\end{astuce}

%------------------------------------------------------------------------------
\subsection{Question 3 : Estimation uniforme}
%------------------------------------------------------------------------------

\textbf{Énoncé :} Montrer qu'il existe une constante $C > 0$ telle que:
\[
\sup_{t \in [0,h]} \E[|X_t - \bar{X}_t|^2] \leq C h^2
\]

\textbf{Correction détaillée :}

\textbf{Étape 1 : Expression de l'erreur}

Pour $t \in [0, h]$:
\[
X_t - \bar{X}_t = \int_0^t [b(X_s) - b(x)] ds + \int_0^t [\sigma(X_s) - \sigma(x)] dW_s
\]

\textbf{Étape 2 : Majoration uniforme}

Par les mêmes arguments que précédemment:
\begin{align*}
\E[|X_t - \bar{X}_t|^2] &\leq 2t \cdot K^2 \int_0^t \E[|X_s - x|^2] ds + 2K^2 \int_0^t \E[|X_s - x|^2] ds \\
&\leq 2(t + 1) K^2 \int_0^t C_1 s \, ds \\
&= 2(t + 1) K^2 C_1 \frac{t^2}{2} \\
&\leq C' t^2 \leq C' h^2
\end{align*}

pour $t \leq h$.

Donc $\sup_{t \in [0,h]} \E[|X_t - \bar{X}_t|^2] \leq C' h^2$.

%------------------------------------------------------------------------------
\subsection{Question 4 : Régularité et borne Lipschitz}
%------------------------------------------------------------------------------

\textbf{Énoncé :} Soit $f : \R \to \R$ une fonction Lipschitz continue. Montrer que:
\[
\E[|f(X_h) - f(\bar{X}_h)|^2] \leq L_f^2 C h^2
\]
où $L_f$ est la constante de Lipschitz de $f$.

\textbf{Correction détaillée :}

Par la condition de Lipschitz de $f$:
\[
|f(X_h) - f(\bar{X}_h)| \leq L_f |X_h - \bar{X}_h|
\]

Donc:
\[
|f(X_h) - f(\bar{X}_h)|^2 \leq L_f^2 |X_h - \bar{X}_h|^2
\]

En prenant l'espérance:
\[
\E[|f(X_h) - f(\bar{X}_h)|^2] \leq L_f^2 \E[|X_h - \bar{X}_h|^2] \leq L_f^2 C h^2
\]

\begin{astuce}
Cette borne est essentielle pour le pricing d'options: si le payoff $f$ est Lipschitz, l'erreur d'approximation Monte Carlo hérite de l'ordre de convergence du schéma d'Euler.
\end{astuce}

%==============================================================================
% PARTIE III : FICHE RÉCAPITULATIVE
%==============================================================================
\newpage
\part{Fiche Récapitulative}

%==============================================================================
\section{Inégalités Fondamentales}
%==============================================================================

\subsection{Inégalités probabilistes}

\begin{enumerate}
    \item \textbf{Inégalité de Cauchy-Schwarz :}
    \[
    |\E[XY]| \leq \sqrt{\E[X^2]\E[Y^2]}
    \]
    \[
    |\Cov(X,Y)| \leq \sigma_X \sigma_Y
    \]
    
    \item \textbf{Inégalité de Jensen :} Si $\varphi$ est convexe:
    \[
    \varphi(\E[X]) \leq \E[\varphi(X)]
    \]
    
    \item \textbf{Inégalité de Markov :} Pour $X \geq 0$ et $a > 0$:
    \[
    \prob(X \geq a) \leq \frac{\E[X]}{a}
    \]
    
    \item \textbf{Inégalité de Bienaymé-Tchebychev :}
    \[
    \prob(|X - \E[X]| \geq \varepsilon) \leq \frac{\Var(X)}{\varepsilon^2}
    \]
    
    \item \textbf{Inégalité de Hölder :} Pour $p, q > 1$ avec $\frac{1}{p} + \frac{1}{q} = 1$:
    \[
    \E[|XY|] \leq \|X\|_p \|Y\|_q = (\E[|X|^p])^{1/p} (\E[|Y|^q])^{1/q}
    \]
\end{enumerate}

\subsection{Inégalités pour les martingales et intégrales stochastiques}

\begin{enumerate}
    \item \textbf{Isométrie d'Itô :}
    \[
    \E\left[ \left( \int_0^T H_s dW_s \right)^2 \right] = \E\left[ \int_0^T H_s^2 ds \right]
    \]
    
    \item \textbf{Inégalité de Burkholder-Davis-Gundy :} Pour $p \geq 1$:
    \[
    \E\left[ \sup_{t \leq T} \left| \int_0^t H_s dW_s \right|^p \right] \leq C_p \E\left[ \left( \int_0^T H_s^2 ds \right)^{p/2} \right]
    \]
    
    \item \textbf{Inégalité de Doob :} Pour une martingale $(M_t)$:
    \[
    \E\left[ \sup_{t \leq T} |M_t|^p \right] \leq \left( \frac{p}{p-1} \right)^p \E[|M_T|^p], \quad p > 1
    \]
\end{enumerate}

\subsection{Inégalités pour les EDS}

\begin{enumerate}
    \item \textbf{Inégalité de Gronwall (forme intégrale) :} Si $u(t) \leq \alpha + \beta \int_0^t u(s) ds$ alors:
    \[
    u(t) \leq \alpha e^{\beta t}
    \]
    
    \item \textbf{Inégalité de Gronwall (forme différentielle) :} Si $u'(t) \leq \beta u(t) + f(t)$ alors:
    \[
    u(t) \leq u(0)e^{\beta t} + \int_0^t e^{\beta(t-s)} f(s) ds
    \]
\end{enumerate}

%==============================================================================
\section{Notions Clés}
%==============================================================================

\subsection{Calcul stochastique}

\begin{enumerate}
    \item \textbf{Formule d'Itô :} Pour $f \in C^2$ et $(X_t)$ processus d'Itô avec $dX_t = \mu_t dt + \sigma_t dW_t$:
    \[
    df(X_t) = f'(X_t) dX_t + \frac{1}{2} f''(X_t) \sigma_t^2 dt
    \]
    \[
    = \left( f'(X_t)\mu_t + \frac{1}{2}f''(X_t)\sigma_t^2 \right) dt + f'(X_t)\sigma_t dW_t
    \]
    
    \item \textbf{Lemme de Stein (intégration par parties gaussienne) :} Pour $Z \sim \mathcal{N}(0,1)$:
    \[
    \E[Zg(Z)] = \E[g'(Z)]
    \]
    
    \item \textbf{Pont brownien :} $Y_t^{0,T} = B_t - \frac{t}{T}(B_T - y)$ est un mouvement brownien conditionné à valoir $y$ en $T$.
    
    \item \textbf{Principe de réflexion :} Pour le mouvement brownien, utilisé pour calculer la loi du supremum.
\end{enumerate}

\subsection{Schémas numériques}

\begin{enumerate}
    \item \textbf{Schéma d'Euler explicite :}
    \[
    \bar{X}_{k+1} = \bar{X}_k + b(\bar{X}_k) h + \sigma(\bar{X}_k) \Delta W_k
    \]
    Convergence forte: $O(\sqrt{h})$, convergence faible: $O(h)$.
    
    \item \textbf{Schéma de Milstein :}
    \[
    \bar{X}_{k+1} = \bar{X}_k + b(\bar{X}_k) h + \sigma(\bar{X}_k) \Delta W_k + \frac{1}{2}\sigma(\bar{X}_k)\sigma'(\bar{X}_k)((\Delta W_k)^2 - h)
    \]
    Convergence forte: $O(h)$.
    
    \item \textbf{Conditions de Lipschitz et croissance linéaire :} Garantissent existence, unicité et stabilité.
\end{enumerate}

\subsection{Monte Carlo et réduction de variance}

\begin{enumerate}
    \item \textbf{Variable de contrôle :} Estimateur $\hat{m}_n^{\lambda} = \frac{1}{n}\sum (X_i - \lambda(Y_i - \mu))$ avec $\lambda^* = \frac{\Cov(X,Y)}{\Var(Y)}$.
    Réduction: facteur $(1 - \rho^2)$.
    
    \item \textbf{Discrépance :} Mesure la régularité de la répartition des points dans $[0,1]^d$.
    \[
    D_n^* = \sup_{u \in [0,1]^d} \left| \frac{1}{n}\sum_{i=1}^n \ind_{[0,u)}(\xi_i) - \prod_{j=1}^d u^j \right|
    \]
    
    \item \textbf{Suite de Halton :} $\xi_n = (\phi_{b_1}(n), \ldots, \phi_{b_d}(n))$ avec discrépance $O\left(\frac{(\log n)^d}{n}\right)$.
    
    \item \textbf{Suite de Van der Corput :} "Miroir" de l'écriture en base $b$ par rapport à la virgule.
\end{enumerate}

%==============================================================================
\section{Astuces et Techniques}
%==============================================================================

\subsection{Techniques de calcul}

\begin{enumerate}
    \item \textbf{Développement de la variance :}
    \[
    \Var(X - aY) = \Var(X) - 2a\Cov(X,Y) + a^2\Var(Y)
    \]
    Minimum en $a^* = \frac{\Cov(X,Y)}{\Var(Y)}$.
    
    \item \textbf{Inégalité triangulaire :} Pour majorer des sommes/intégrales:
    \[
    |a + b| \leq |a| + |b|, \quad |a + b|^2 \leq 2(|a|^2 + |b|^2)
    \]
    
    \item \textbf{Inégalité de Cauchy-Schwarz pour les intégrales :}
    \[
    \left( \int_0^T f(t) dt \right)^2 \leq T \int_0^T f(t)^2 dt
    \]
    
    \item \textbf{Utilisation du Gronwall :} Pour les équations d'erreur du type:
    \[
    \E[\sup_{s \leq t} |e_s|^2] \leq C_1 \int_0^t \E[\sup_{u \leq s} |e_u|^2] ds + C_2 h
    \]
\end{enumerate}

\subsection{Astuces pour les EDS}

\begin{enumerate}
    \item \textbf{Positivité :} Pour montrer qu'un processus reste positif, utiliser la représentation exponentielle (formule d'Itô avec $f(x) = \log x$ ou $f(x) = e^x$).
    
    \item \textbf{Changement de variable :} Pour simplifier une EDS, chercher $Y_t = f(X_t)$ tel que $Y$ suive une EDS plus simple (ex: linéaire).
    
    \item \textbf{Encadrement des erreurs :} Séparer la partie drift et la partie martingale, traiter séparément.
\end{enumerate}

\subsection{Astuces pour Quasi-Monte Carlo}

\begin{enumerate}
    \item \textbf{Structure de la suite de Hammersley :} Fixer la dernière coordonnée comme $\frac{k}{n}$ réduit la discrépance d'un facteur $\log n$.
    
    \item \textbf{Fonction affine par morceaux :} Le supremum est atteint aux points de non-dérivabilité ou aux bords.
    
    \item \textbf{Bases premières :} Pour Halton, utiliser des bases deux à deux premières entre elles pour éviter les corrélations.
\end{enumerate}

%==============================================================================
\section{Formules à Retenir}
%==============================================================================

\begin{tcolorbox}[colback=yellow!10!white,colframe=orange!75!black,title=Formules Essentielles]

\textbf{Variance de la régression :}
\[
\min_a \Var(X - aY) = \Var(X)(1 - \rho^2)
\]

\textbf{Convergence forte Euler :}
\[
\E\left[\sup_{t \leq T} |X_t - \bar{X}_t^n|^2\right] = O(h) = O(1/n)
\]

\textbf{Discrépance de Halton :}
\[
D_n^* = O\left(\frac{(\log n)^d}{n}\right)
\]

\textbf{Loi du sup du pont brownien :} Pour $s \geq \max(0, y)$:
\[
\prob\left(\sup_{t \in [0,T]} Y_t^{0,T} > s\right) = e^{-\frac{2s(s-y)}{T}}
\]

\textbf{Lemme de Stein :}
\[
\E[Zg(Z)] = \E[g'(Z)] \quad \text{pour } Z \sim \mathcal{N}(0,1)
\]

\end{tcolorbox}

\end{document}
